\section{Architecture of Certificates}
\label{sec:certificates}

\subsection{Strong public certificates}

The public seam distinguishes transformation equations from invertibility.
The equation-only structures are \code{LeftTransformEquation} and
\code{TwoSidedTransformEquation}. A strong result instead uses:

\begin{lstlisting}[style=leanstyle,caption={Explicit inverse certificate seam.}]
structure MatrixInverseCertificate (U : Matrix n n R) where
  inverse   : Matrix n n R
  left_inv  : inverse * U = 1
  right_inv : U * inverse = 1

structure LeftCertificate (A : Matrix m n R) where
  U        : Matrix m m R
  U_cert   : MatrixInverseCertificate U
  result   : Matrix m n R
  equation : U * A = result
\end{lstlisting}

\code{TwoSidedCertificate} analogously stores \(U,U^{-1},V,V^{-1}\), both
inverse identities on each side, and \(UAV=S\). Matrix-unit and
determinant-based unimodularity views are derived from this data rather than
stored independently.

\subsection{Constructive composition}

Inverse certificates compose with the inverse product in reverse order. They
also support symmetry, transpose, and explicit reindexing. Elementary swap,
row/column addition, unit scaling, and the \(2\times2\) Bezout primitive each
provide their inverse matrix directly. Recursive diagonal and leading-block
lifts transport those inverses structurally. Consequently, the algorithmic
path never calls \code{Matrix.inv}, \code{nonsing_inv}, or a determinant
argument to obtain inverse data.

The HNF transform stores \(U\), its certificate, and \(UA=H\). The SNF
transform stores both certified sides and \(UAV=S\). Composition preserves the
same strong invariant. Public \code{HNFResultFin} and \code{SNFResultFin}
therefore expose the exact data produced by the recursive kernels.

\subsection{Reversible logs}

The internal reversible-log layer makes small-step execution equally explicit.
Every step stores forward and backward matrices and both inverse proofs. Replay
accumulates four matrices:
\[
  U,\qquad U^{-1},\qquad V,\qquad V^{-1}.
\]
The library proves both identities for each accumulator, the replay equation
\(UAV=\operatorname{replay}(A)\), and the ordering laws for log append. A
strong two-sided certificate is obtained directly from the log. Raw,
potentially noninvertible operations remain useful for testing execution, but
they produce only equation-level evidence.

\subsection{Index transport}

For any two \code{MatrixIndexing} values the library constructs explicit row
and column permutation matrices, their inverse certificates, and the equation
relating the two reindexed inputs. This data is then composed with result
certificates. For HNF it yields left equivalence between outputs. For SNF it
yields a two-sided equivalence between canonical \code{Fin} matrices, which
the Smith uniqueness theorem upgrades to equality.

\subsection{Untrusted external certificates}

The stable JSON schema \code{normalforms.cert/v1} represents every integer by
the canonical grammar \(0\) or \code{-?[1-9][0-9]*}. A parser produces raw HNF
or SNF data. The pure checkers \code{checkHNF} and \code{checkSNF} take the
expected matrix as a separate caller argument; they reject dimension errors,
an input mismatch, the transformation equation, either inverse direction, or
the normal-form predicate independently. Checked data is only typed data. A
strong result follows from \code{checkHNF\_sound} or
\code{checkSNF\_sound} and an equation showing checker success.
This proof-search/proof-checking separation follows the skeptical integration
principle of Harrison and Th\'ery \citep{harrisonThery1998}.

The \code{normalforms-check emit-lean} command renders a module importing the
caller's matrix, a typed raw certificate literal, and separate Boolean
obligations for dimensions, input identity, transformation, inverses, and
normal form. Each obligation is closed with \code{decide_cbv}; the final
declaration uses only checker soundness. At Lean 4.32.0 this tactic performs
call-by-value reduction and does not rely on the code generator
\citep{lean4320decideCbv}. The strict declaration audit permits
\code{propext} and \code{Quot.sound}, but rejects
\code{Classical.choice}, native reduction, and compiler trust.

FLINT is isolated outside this trusted path. Against the pinned FLINT 3.6.0
integer-matrix interface \citep{flint360fmpzmat}, the optional C computation uses
\code{fmpz\_mat\_hnf\_transform}, \code{fmpz\_mat\_snf\_transform}, and
\code{fmpz\_mat\_inv}; it accepts inverse output only when the common
denominator is \(1\) or \(-1\). The process worker and Lean FFI invoke the same
strict decimal-protocol function with standard and memory streams,
respectively. A Python standard-library wrapper validates the protocol and
emits JSON; the Lean FFI layer independently parses it back through the same
raw schema and pure checker. Neither route returns a proof.

Seeded differential tests compare only the normalized nonzero diagonal
factors, because \(U\) and \(V\) are not unique. Successful generation,
differential agreement, and native checker acceptance are run reports, not
theorems: the emitted Lean module must still compile. The fixed version 1
benchmark records one warmup and seven measurements per stage. On the captured
host, median process and FFI generation costs are \(0.03\) and \(0.06\) seconds,
native checking is \(0.02\) seconds, and strict kernel checking is \(7.04\)
seconds. The raw JSON/CSV, certificate sizes, and hardware fingerprint are the
authoritative data; no absolute time is a portable CI threshold.
